
\section{Principes POO mobilisés}
\begin{itemize}[leftmargin=*]
  \item \textbf{Encapsulation :} Ce principe est rigoureusement appliqué. Les attributs des classes (par exemple, \texttt{solde} dans \texttt{Compte}, \texttt{salaire} dans \texttt{Employe}) sont privés. L'accès et la modification de ces données sont contrôlés exclusivement par des méthodes publiques (accesseurs et mutateurs), garantissant ainsi la cohérence et l'intégrité de l'état des objets. 

  \item \textbf{Héritage et Hiérarchie de Classes :} L'héritage est au cœur de la conception. Il est utilisé pour créer une hiérarchie de classes claire : \texttt{Compte} est la base pour \texttt{CompteCourant} et \texttt{CompteEpargne}. De même, \texttt{Utilisateur} est la base pour \texttt{Client} et \texttt{Employe}, cette dernière servant à son tour de base pour des rôles spécifiques comme \texttt{Manager}, \texttt{Caissier}, et \texttt{AdminIT}. 

  \item \textbf{Héritage Virtuel :} Le projet utilise de manière avancée l'héritage virtuel pour résoudre le "problème du diamant". La classe \texttt{EmployeClient} hérite à la fois de \texttt{Employe} et \texttt{Client}, qui héritent eux-mêmes virtuellement de \texttt{Utilisateur}. Cela garantit qu'un objet \texttt{EmployeClient} ne contienne qu'une seule et unique sous-partie de la classe \texttt{Utilisateur}, évitant ainsi l'ambiguïté et la duplication de données. 

  \item \textbf{Abstraction :} L'abstraction est réalisée à deux niveaux. Au niveau des classes, des classes de base abstraites comme \texttt{Utilisateur} et \texttt{Compte} définissent des interfaces communes avec des fonctions virtuelles pures (ex: \texttt{afficherProfil()}, \texttt{retirer()}), forçant les classes dérivées à fournir une implémentation concrète. Au niveau de l'architecture, le système est abstrait en couches logiques (présentation, logique métier, accès aux données) grâce aux classes \texttt{InterfaceUtilisateur}, aux Managers, et au \texttt{BDManager}. 

  \item \textbf{Polymorphisme :} Le polymorphisme est intensivement utilisé pour manipuler des objets de types différents via une interface commune. Les conteneurs comme \texttt{std::vector<std::unique\_ptr<Utilisateur>>} dans \texttt{UtilisateurManager} permettent de stocker et de traiter des \texttt{Client}s, \texttt{Manager}s, etc., de manière uniforme. L'appel à des méthodes virtuelles (\texttt{virtual}) redéfinies (\texttt{override}) assure que le comportement correct est exécuté à l'exécution (liaison tardive). 

  \item \textbf{Composition sur Agrégation :} Plutôt qu'une simple agrégation, le projet privilégie la composition pour gérer les relations de possession. Les classes \texttt{UtilisateurManager} et \texttt{CompteManager} possèdent et gèrent le cycle de vie des utilisateurs et des comptes via des pointeurs intelligents (\texttt{std::unique\_ptr}). Cela garantit que la mémoire est gérée automatiquement, prévenant les fuites et clarifiant la sémantique de possession. 

  \item \textbf{Patrons de Conception (Design Patterns) :} L'architecture emploie des patrons de conception éprouvés. Le \textbf{Singleton}, utilisé pour la classe \texttt{BDManager}, assure qu'il n'existe qu'une seule instance de connexion à la base de données dans toute l'application. La séparation en Managers et Services s'apparente au principe de \textbf{Séparation des Préoccupations (SoC)}, un fondement des architectures logicielles robustes. 

  \item \textbf{Exceptions personnalisées :} Pour une gestion des erreurs claire et sémantique, le projet définit des exceptions spécifiques (\texttt{CompteIntrouvable}, \texttt{SoldeInsuffisant}). Cela permet de distinguer les différents types d'erreurs d'exécution et de mettre en place des mécanismes de récupération plus fins qu'avec des codes d'erreur ou des exceptions génériques.
\end{itemize}

\section{Conclusion du chapitre}
La conception de ce projet bancaire dépasse une simple application des principes orientés objet ; elle témoigne d'une architecture logicielle réfléchie et moderne. En adoptant une \textbf{séparation claire des préoccupations}, le système isole la logique de présentation (\texttt{InterfaceUtilisateur}), la logique métier (\texttt{UtilisateurManager}, \texttt{CompteManager}) et l'accès aux données (\texttt{BDManager}). Cette modularité, renforcée par l'usage de patrons de conception comme le Singleton, n'est pas un simple exercice académique : elle confère au projet une maintenabilité et une évolutivité concrètes.

L'utilisation de fonctionnalités C++ modernes, notamment les \textbf{pointeurs intelligents (\texttt{std::unique\_ptr})} pour la gestion automatique de la mémoire et \textbf{l'héritage virtuel} pour résoudre des cas de figure complexes, démontre une maîtrise qui va au-delà des bases de la POO. Ces choix techniques fiabilisent l'application en prévenant les fuites de mémoire et en assurant une hiérarchie de classes cohérente.

En conclusion, l'architecture mise en place ne se contente pas de structurer le code. Elle le rend plus sûr, plus facile à maintenir et prêt pour des évolutions futures. L'ajout de nouveaux types de comptes, de rôles d'employés ou même le passage à une autre technologie d'interface utilisateur pourrait se faire avec un impact minimal sur le cœur du système, prouvant ainsi la robustesse et la qualité de la conception initiale.
